% Diseño, implementación y evidencia de desplegabilidad, comprobabilidad,
% integrabilidad, usabilidad y portabilidad sobre la arquitectura vigente.
% Responsables: Diego Coronado (RNF-15, RNF-18) y Daniel Cristancho (RNF-16,
% RNF-12, RNF-14). Lo incluye DescripcionArquitecturaSoftware.tex con \input.
% Las cifras salen de INTEGRABILIDAD-RNF16.md, USABILIDAD-PORTABILIDAD.md,
% App/infra/DESPLIEGUE.md y de las ejecuciones citadas.

\section{Desplegabilidad, comprobabilidad, integrabilidad, usabilidad y portabilidad}
\label{sec:atributos-equipo}

Esta sección completa el diseño de los atributos de calidad del prototipo con los que no se
documentan en las secciones \ref{sec:atributos-vigentes} y \ref{sec:seguridad-vigente}:
\textbf{desplegabilidad} (RNF-15) y \textbf{comprobabilidad} (RNF-18), a cargo de Diego Coronado;
\textbf{integrabilidad} (RNF-16), \textbf{usabilidad} (RNF-12) y \textbf{portabilidad} (RNF-14), a
cargo de Daniel Cristancho. Igual que en las secciones anteriores, para cada atributo se da el
escenario de calidad, las tácticas con el mecanismo que las implementa, la evidencia y lo que
todavía no está acreditado. A diferencia de la sección \ref{sec:tacticas}, todo se refiere a la
arquitectura vigente (ADR-16).

% ------------------------------------------------------------------
\subsection{Desplegabilidad}
\label{sec:desplegabilidad-vigente}

\begin{center}
\begin{tabular}{|>{\raggedright\arraybackslash}p{3.2cm}|>{\raggedright\arraybackslash}p{10.4cm}|}
\hline
\multicolumn{2}{|l|}{\textbf{Escenario de calidad · Desplegabilidad (RNF-15)}}\\
\hline
Fuente & Un integrante del equipo que fusiona un cambio en \texttt{main}.\\
\hline
Estímulo & Una versión nueva de cualquiera de los programas.\\
\hline
Artefacto & Las ocho imágenes del sistema y el servidor de presentación.\\
\hline
Entorno & Operación normal del servidor, con datos.\\
\hline
Respuesta & GitHub Actions compila, publica, despliega y comprueba sin intervención manual. Si la
versión nueva no queda sana, el servidor vuelve solo a la anterior.\\
\hline
Medida & Ventana de despliegue $\leq$ 10 min, reversión $\leq$ 5 min y 0 solicitudes fallidas
durante el despliegue.\\
\hline
\end{tabular}
\end{center}

\begin{center}
\begin{longtable}{|>{\raggedright\arraybackslash}p{2.8cm}|>{\raggedright\arraybackslash}p{6.7cm}|>{\raggedright\arraybackslash}p{3.8cm}|}
\caption{Tácticas de desplegabilidad en la arquitectura vigente.}\label{tab:desplegabilidad-vigente}\\
\hline
\textbf{Táctica} & \textbf{Mecanismo} & \textbf{Dónde}\\
\hline
\endfirsthead
\hline
\textbf{Táctica} & \textbf{Mecanismo} & \textbf{Dónde}\\
\hline
\endhead
Automatizar el despliegue &
Cada merge a \texttt{main} pasa por el CI y, si sale verde, por el CD: compila las ocho imágenes
para amd64 y arm64, las publica en GHCR y comprueba sus digests antes de desplegar. &
\texttt{.github/\allowbreak workflows/\allowbreak cd.yml}\\
\hline
Versiones identificables &
Se despliega siempre una etiqueta inmutable \texttt{sha-<commit>}, y el servidor anota la versión
activa. Gateway y Caddy exponen \texttt{/version}. &
\texttt{desplegar.sh}, \texttt{.version-\allowbreak desplegada}\\
\hline
Reversión automática &
Si la versión nueva no llega a estar sana, \texttt{desplegar.sh} levanta la anterior y el trabajo
del CD falla en rojo. Las migraciones solo van hacia delante: revertir cambia imágenes, no datos. &
\texttt{desplegar.sh}, \texttt{CORTE-Y-\allowbreak REVERSION.md}\\
\hline
Configuración por entorno &
El mismo artefacto sirve en local, en la demo y en producción; cambian el \texttt{.env}, el dominio
y las claves JWT propias del servidor. &
\texttt{preparar-\allowbreak servidor.sh}, \texttt{docker-compose.\allowbreak produccion.yml}\\
\hline
Acceso controlado al servidor &
El CD entra por SSH con una llave exclusiva y la huella del servidor fijada en un secreto, en vez
de aceptar la que diga el servidor al conectar. &
Secretos \texttt{PRODUCCION\_SSH\_*}\\
\hline
\end{longtable}
\end{center}

\begin{sloppypar}\noindent\textbf{Evidencia.} El servidor de presentación corre en Google Cloud (\texttt{e2-standard-2},
Ubuntu 24.04) con HTTPS de Let's Encrypt en
\href{https://hexacore-javeriana.duckdns.org}{\texttt{hexacore-\allowbreak javeriana.\allowbreak duckdns.org}} y
\texttt{admin.\allowbreak hexacore-\allowbreak javeriana.\allowbreak duckdns.org}
(\href{https://github.com/Emperador24/HEXACORE/blob/main/App/infra/DESPLIEGUE.md}{DESPLIEGUE}). En
la ejecución
\href{https://github.com/Emperador24/HEXACORE/actions/runs/43987694799}{43987694799} del CD, el
trabajo «Desplegar en producción» entró por SSH, desplegó \texttt{sha-9815d8e} y comprobó que
\texttt{/api/v1/salud} respondiera, en \textbf{62 s}: dentro de los 10 min de RNF-15.\end{sloppypar}

\paragraph{No acreditado.} No se ha provocado en vivo una versión que falle para medir la reversión
(5 min), ni se han contado las solicitudes fallidas durante un despliegue: el compose de
producción tiene una réplica por programa, así que al recrear un contenedor hay una ventana sin
servicio que no está medida.

% ------------------------------------------------------------------
\subsection{Comprobabilidad}
\label{sec:comprobabilidad-vigente}

\begin{center}
\begin{tabular}{|>{\raggedright\arraybackslash}p{3.2cm}|>{\raggedright\arraybackslash}p{10.4cm}|}
\hline
\multicolumn{2}{|l|}{\textbf{Escenario de calidad · Comprobabilidad (RNF-18)}}\\
\hline
Fuente & Un integrante que cambia un módulo, o el CI en cada pull request.\\
\hline
Estímulo & Ejecutar las pruebas de un módulo.\\
\hline
Artefacto & Los módulos del Núcleo, Boletería y Trabajador.\\
\hline
Entorno & Desarrollo y CI, sin levantar los demás programas.\\
\hline
Respuesta & Cada módulo se prueba aislado: base, broker, pasarela y correo se reemplazan por dobles.\\
\hline
Medida & Cobertura de la lógica de dominio $\geq$ 70\% y 100\% de las dependencias externas
sustituibles por dobles.\\
\hline
\end{tabular}
\end{center}

\begin{center}
\begin{longtable}{|>{\raggedright\arraybackslash}p{2.8cm}|>{\raggedright\arraybackslash}p{6.7cm}|>{\raggedright\arraybackslash}p{3.8cm}|}
\caption{Tácticas de comprobabilidad en la arquitectura vigente.}\label{tab:comprobabilidad-vigente}\\
\hline
\textbf{Táctica} & \textbf{Mecanismo} & \textbf{Dónde}\\
\hline
\endfirsthead
\hline
\textbf{Táctica} & \textbf{Mecanismo} & \textbf{Dónde}\\
\hline
\endhead
Abstraer las dependencias externas &
El negocio llega a lo externo por interfaces propias (\texttt{ProcesadorPagos}, correo, broker), que
las pruebas reemplazan por dobles. La pasarela y el correo tienen además simuladores HTTP. &
\texttt{modulos/\allowbreak finanzas/\allowbreak pagos}, \texttt{App/infra/\allowbreak pasarela-pagos-\allowbreak simulada}\\
\hline
Separar niveles de prueba &
Unitarias (Jest) sin infraestructura; pruebas HTTP contra los programas reales; contratos de
mensajería con RabbitMQ real; pruebas de imagen. Cada nivel tiene su configuración. &
\texttt{jest-\allowbreak unitarias.json}, \texttt{test/\allowbreak jest-e2e.json}\\
\hline
Entornos de prueba exclusivos &
Las pruebas que necesitan base o broker levantan los suyos, con nombres propios, sin tocar el
entorno de desarrollo. &
\texttt{test/README.md}\\
\hline
Hacer cumplir las fronteras &
Una comprobación del CI impide que un módulo importe archivos internos de otro: solo su API
pública. Así cada módulo se puede probar y reemplazar por separado. &
\texttt{test/\allowbreak comprobar-\allowbreak configuracion.py}\\
\hline
Ejecutar todo en cada cambio &
El CI corre en cada pull request las suites del backend, los portales, la app móvil y el gateway. &
\texttt{.github/\allowbreak workflows/\allowbreak ci.yml}\\
\hline
\end{longtable}
\end{center}

\paragraph{Evidencia.} En el corte del 17/11/2026 pasan \textbf{1\,162 pruebas unitarias} del
backend (39 omitidas), 90 de la app móvil y las suites de los portales; el CI del PR~66 salió
verde en todos sus trabajos. Midiendo solo las pruebas unitarias
(\texttt{npm run test:cov}), la cobertura por línea es \textbf{68,7\%} sobre los archivos que
cargan, y \textbf{62,4\%} sobre los servicios y el dominio de los módulos.

\paragraph{No acreditado.} Esas cifras no alcanzan el 70\% de RNF-18 contando solo unitarias.
\texttt{npm run test:cov:total} une por línea la cobertura de las unitarias y de las pruebas HTTP,
que ejecutan los programas reales, pero ese total no se ha publicado para esta revisión, ni el
reporte de cobertura forma parte de GitHub Actions. Tampoco hay una matriz que acredite el 100\%
de dependencias sustituibles.

% ------------------------------------------------------------------
\subsection{Integrabilidad}
\label{sec:integrabilidad-vigente}

\begin{center}
\begin{tabular}{|>{\raggedright\arraybackslash}p{3.2cm}|>{\raggedright\arraybackslash}p{10.4cm}|}
\hline
\multicolumn{2}{|l|}{\textbf{Escenario de calidad · Integrabilidad (RNF-16)}}\\
\hline
Fuente & La organización, que decide cambiar de proveedor de pagos.\\
\hline
Estímulo & Reemplazar la pasarela simulada por Stripe, o por otro proveedor.\\
\hline
Artefacto & El componente Pagos y los módulos que cobran: Entradas, Alimentos y Parqueaderos.\\
\hline
Entorno & Sistema desplegado, sin recompilar.\\
\hline
Respuesta & Se cambia una variable y los cobros, rechazos y reembolsos siguen funcionando.\\
\hline
Medida & 1 componente modificado (el adaptador) y 0 cambios en Entradas.\\
\hline
\end{tabular}
\end{center}

\begin{center}
\begin{longtable}{|>{\raggedright\arraybackslash}p{2.8cm}|>{\raggedright\arraybackslash}p{6.7cm}|>{\raggedright\arraybackslash}p{3.8cm}|}
\caption{Tácticas de integrabilidad en la arquitectura vigente.}\label{tab:integrabilidad-vigente}\\
\hline
\textbf{Táctica} & \textbf{Mecanismo} & \textbf{Dónde}\\
\hline
\endfirsthead
\hline
\textbf{Táctica} & \textbf{Mecanismo} & \textbf{Dónde}\\
\hline
\endhead
Encapsular (interfaz propia) &
El negocio cobra y reembolsa con \texttt{ProcesadorPagos} y siempre recibe uno de tres resultados:
\texttt{APROBADO}, \texttt{RECHAZADO} o \texttt{INDETERMINADO}, iguales para cualquier proveedor. &
\texttt{modulos/\allowbreak finanzas/\allowbreak pagos}\\
\hline
Adaptar &
Un adaptador por proveedor: \texttt{PasarelaHttp} (simulada) y \texttt{PasarelaStripe}
(PaymentIntents y Refunds). Traducen el protocolo de cada uno al contrato propio. &
\texttt{pasarela-stripe.\allowbreak service.ts}, \texttt{pasarela-http.\allowbreak service.ts}\\
\hline
Vincular en la configuración &
El único lugar que sabe qué pasarela hay es \texttt{crearProcesadorPagos}, que la elige por
\texttt{PAGOS\_PROVEEDOR}. Cambiar de proveedor es cambiar una variable y reiniciar. &
\texttt{pagos.module.ts}\\
\hline
Restringir dependencias &
Los módulos de negocio llegan a Pagos solo por su API pública (\texttt{index.ts}); una prueba
comprueba que ninguno nombra a Stripe, a un adaptador o a la variable. &
\texttt{rnf16-cambio-\allowbreak pasarela.spec.ts}\\
\hline
\end{longtable}
\end{center}

\paragraph{Protocolos.} Internos: HTTP/JSON por el API Gateway (\texttt{/api/v1}, JWT RS256
verificado en el gateway), AMQP con exchanges \emph{topic} entre módulos (bandeja de salida
transaccional) y SSE para los paneles en vivo. Externo: Stripe, con la tarjeta tokenizada en el
navegador por Stripe.js (la tarjeta nunca llega a HEXACORE), cobro por \texttt{POST
/v1/payment\_intents} con \texttt{Idempotency-Key}, reembolso por \texttt{POST /v1/refunds} y
confirmación tardía por un webhook firmado con HMAC-SHA256.

\paragraph{Evidencia.} La prueba \texttt{rnf16-cambio-pasarela.spec.ts} (21 casos, en el CI) pasa la
misma tabla de solicitudes por los dos adaptadores y obtiene los mismos resultados de negocio. En el
sandbox de Stripe se cobró, se rechazó y se reembolsó de verdad por el camino completo. El cambio
de pasarela se aplicó en \textbf{8 s}, sin recompilar y con \textbf{0 archivos de código
modificados}
(\href{https://github.com/Emperador24/HEXACORE/blob/main/Documentation/Work/INTEGRABILIDAD-RNF16.md}{INTEGRABILIDAD-RNF16}).
La prueba real encontró y corrigió que el reembolso de un cobro de Stripe se rechazaba dentro de
Pagos.

\paragraph{Límites.} Un proveedor nuevo también necesita su formulario de tarjeta en los portales,
porque la tokenización ocurre en el navegador; RNF-16 mide el backend. Stripe en modo de prueba no
permite forzar un reembolso rechazado.

% ------------------------------------------------------------------
\subsection{Usabilidad y arquitectura de información}
\label{sec:usabilidad-vigente}

\begin{center}
\begin{tabular}{|>{\raggedright\arraybackslash}p{3.2cm}|>{\raggedright\arraybackslash}p{10.4cm}|}
\hline
\multicolumn{2}{|l|}{\textbf{Escenario de calidad · Usabilidad (RNF-12)}}\\
\hline
Fuente & Un cliente que nunca ha usado HEXACORE.\\
\hline
Estímulo & Quiere comprar una entrada (CU-001).\\
\hline
Artefacto & Portal web de clientes y app móvil.\\
\hline
Entorno & Uso normal, sin capacitación previa.\\
\hline
Respuesta & Completa la compra y ve su QR sin ayuda.\\
\hline
Medida & $\geq$ 80\% de éxito en una prueba con 5 usuarios, en $\leq$ 5 pasos desde el evento hasta
la confirmación.\\
\hline
\end{tabular}
\end{center}

Cada grupo de pantallas sigue un patrón de arquitectura de información según la tarea que
resuelve:

\begin{center}
\begin{longtable}{|>{\raggedright\arraybackslash}p{3.2cm}|>{\raggedright\arraybackslash}p{2.5cm}|>{\raggedright\arraybackslash}p{7.6cm}|}
\caption{Patrones de arquitectura de información por grupo de pantallas.}\label{tab:arquitectura-informacion}\\
\hline
\textbf{Grupo de pantallas} & \textbf{Patrón} & \textbf{Cómo se ve en el prototipo}\\
\hline
\endfirsthead
\hline
\textbf{Grupo de pantallas} & \textbf{Patrón} & \textbf{Cómo se ve en el prototipo}\\
\hline
\endhead
Compra de entradas (CU-001) y pagos de pedidos y parqueadero & Secuencial &
Una sola ruta hacia delante: evento $\rightarrow$ localidad y cantidad $\rightarrow$ pago
$\rightarrow$ confirmación con QR. La cola virtual y la reserva ocurren dentro del mismo paso.\\
\hline
Cartelera de eventos (CU-005) & Matricial &
Los eventos se cruzan por varios criterios combinables a la vez: búsqueda por texto, categoría,
fecha y ciudad.\\
\hline
Registro, verificación y recuperación de cuenta (CU-027) & Secuencial &
Pasos encadenados por el correo: registrarse, verificar con el enlace, iniciar sesión; o pedir
el enlace y fijar la contraseña nueva.\\
\hline
Navegación del cliente (web y app) & Jerárquico plano (centro y radios) &
Pocas secciones de primer nivel siempre visibles (en la app, la barra inferior: Eventos,
Parqueadero, Pedidos, Reventa y Mis pedidos), y desde cada una su lista y su detalle.\\
\hline
Portal administrativo & Jerárquico, filtrado por rol &
Menú lateral en tres grupos (General, Operación y Administración) $\rightarrow$ sección
$\rightarrow$ lista $\rightarrow$ detalle o formulario nuevo (por ejemplo,
\texttt{eventos}, \texttt{eventos/nuevo}, \texttt{eventos/:id}). Cada cuenta ve solo las
secciones de sus roles; la barra superior muestra la ubicación (sección y página).\\
\hline
Gestión de personal (CU-007 a CU-009) & Jerárquico con pestañas &
Dentro de Personal, pestañas para empleados, asignaciones, cambios de turno y asistencia: cuatro
vistas de la misma información, al mismo nivel.\\
\hline
Paneles en vivo: monitoreo, parqueaderos y emergencias (CU-019, CU-025, CU-010) & Tablero &
Una sola pantalla con las zonas y sus indicadores, que se actualiza sola (SSE); el detalle y las
acciones se abren desde ahí sin salir del panel.\\
\hline
App del personal & Por audiencia (rol y puesto) &
La pestaña operativa depende del puesto del turno: Entrada ve validar entradas; Parqueadero, el
parqueadero; Restaurante, menú y pedidos. Todos ven además turnos, asistencia, incidentes y
emergencia.\\
\hline
\end{longtable}
\end{center}

\paragraph{Evidencia.} Por diseño, la compra toma \textbf{4 pasos} en la app y \textbf{5} en la web
desde el detalle del evento hasta ver el QR, dentro del límite de RNF-12; el código promocional es
opcional y suma uno. Las demostraciones automatizadas recorren esos flujos en los portales reales
(\href{https://github.com/Emperador24/HEXACORE/blob/main/Documentation/Work/USABILIDAD-PORTABILIDAD.md}{USABILIDAD-PORTABILIDAD}).

\paragraph{No acreditado.} La prueba con cinco usuarios nuevos no se ha hecho: el porcentaje de
éxito de RNF-12 es desconocido. El protocolo (consigna, criterio de éxito y plantilla) está
escrito y listo para ejecutarse.

% ------------------------------------------------------------------
\subsection{Portabilidad}
\label{sec:portabilidad-vigente}

\begin{center}
\begin{tabular}{|>{\raggedright\arraybackslash}p{3.2cm}|>{\raggedright\arraybackslash}p{10.4cm}|}
\hline
\multicolumn{2}{|l|}{\textbf{Escenario de calidad · Portabilidad (RNF-14)}}\\
\hline
Fuente & Un cliente.\\
\hline
Estímulo & Usa una funcionalidad de su rol desde el navegador o desde el teléfono.\\
\hline
Artefacto & Portal web de clientes, app móvil y la API del gateway.\\
\hline
Entorno & Operación normal.\\
\hline
Respuesta & La misma funcionalidad está en las dos interfaces, sobre el mismo contrato.\\
\hline
Medida & 100\% de los CU del rol Cliente en ambas interfaces, con 0 reglas de negocio duplicadas en
el cliente.\\
\hline
\end{tabular}
\end{center}

\paragraph{Diseño.} Las dos interfaces son clientes delgados del mismo API Gateway
(\texttt{/api/v1}): precios, totales, descuentos, cupones y reembolsos los calcula el servidor, y los
clientes solo validan la forma de lo que envían (largo de un texto, formato de una placa). La
política de cancelación llega en la cotización y el cliente confirma con ese monto; si cambió, el
servidor rechaza.

\paragraph{Evidencia.} \textbf{11 de 11} CU con parte interactiva del Cliente (CU-001 a CU-006,
CU-011, CU-013/014, CU-021, CU-022 y CU-027) están en la web y en la app sobre los mismos
endpoints; CU-024 no tiene pantalla de Cliente en ninguna, por diseño del caso de uso. Antes de esta
entrega la matriz daba 6 de 11 y el descuento se calculaba en el teléfono; se reemplazaron las
pantallas simuladas de la app por las reales. Lo comprueban \texttt{boleteria\_test.dart}, la suite
completa de la app (90 pruebas) y \texttt{entradas.service.spec.ts} del portal.

\paragraph{Límites.} La app se probó en emulador Android: falta un teléfono físico contra la URL
HTTPS de la demostración. La app usa los tokens de la pasarela simulada, que en modo de prueba se
traducen a tarjetas de prueba de Stripe; no tiene formulario de tarjeta propio de Stripe.
